\chapter{Определения системы искусственного интеллекта (СИИ)}

\textit{Определение 1}: Под СИИ понимается аппаратно-программный комплекс, обладающий способностью к накоплению и корректировке знаний на основе активного восприятия информации о мире и обобщённого опыта; к целенаправленному поведению на основе логического вывода и алгоритмов различной степени неопределённости представления и управления сложных систем.

\textit{Определение 2}:  СИИ = ЗНАНИЯ + СТРАТЕГИЯ ОБРАБОТКИ ЗНАНИЙ (в отличие от обычной программы: ПРОГРАММА = АЛГОРИТМ + ДАННЫЕ).

\textit{Мое гибридное определение}: СИИ — это аппаратно-программный комплекс, работающий по схеме «знания + стратегия обработки знаний», который накапливает и корректирует знания, целенаправленно обрабатывает их, решая задачи, традиционно считающиеся интеллектуальными.

\chapter{Признаки ИИ у разрабатываемой рекомендательной системы}
\begin{enumerate}
\item Работа со знаниями о данных, а не с сырыми данными ("СИИ = ЗНАНИЯ + СТРАТЕГИЯ ОБРАБОТКИ ЗНАНИЙ"): критерии сравнения, выделенные атрибуты, таксономия, формулы обработки и оценивания данных.
\item Способность к "накоплению и корректировке знаний на основе активного восприятия информации о мире и обобщённого опыта": накопление и корректировка знаений при проставлении пользователем лайков и дизлайков, что влияло на степень схожести книг
\item Использует знания «для постановки и достижения цели (целенаправленного поведения), адаптации к изменениям среды и внутреннего состояния»: целью является сортировка книг по релевантности, в реализованной программе коэффиценты релевантности адаптируются под изменения внешней среды (проставление лайков и дизлайков)
\end{enumerate}


\chapter{Сравнительный анализ методов формирования рекомендаций}

Для сравнения методов формирования рекомендаций были выбраны следующие критерии:

\begin{enumerate}
	\item \textit{Проблема холодного старта} --- сложность формирования рекомендаций для нового пользователя из-за отсутствия информации о его предпочтениях.
	
	\item \textit{Проблема белой вороны} --- сложность формирования рекомендаций для пользователей с нетипичными предпочтениями, для которых трудно найти похожих пользователей.
	
	\item \textit{Проблема устаревания предпочтений} --- снижение актуальности рекомендаций вследствие изменения интересов пользователя с течением времени.
\end{enumerate}

\begin{longtable}{|p{0.20\textwidth}|p{0.23\textwidth}|p{0.23\textwidth}|p{0.23\textwidth}|}
	\caption{Сравнительный анализ методов формирования рекомендаций}
	\label{tab:methods}\\ \hline
	
	\textbf{Критерий сравнения} &
	\textbf{Коллаборативная фильтрация (User-Based)} &
	\textbf{Коллаборативная фильтрация (Item-Based)} &
	\textbf{Контентная фильтрация} \\ \hline
	\endfirsthead
	
	\hline
	\textbf{Критерий сравнения} &
	\textbf{Коллаборативная фильтрация (User-Based)} &
	\textbf{Коллаборативная фильтрация (Item-Based)} &
	\textbf{Контентная фильтрация} \\ \hline
	\endhead
	
	Краткое описание &
	Поиск пользователей со схожими предпочтениями &
	Рекомендация объектов, похожих на ранее оценённые &
	Сопоставление характеристик объектов с предпочтениями пользователя \\ \hline
	
	Проблема холодного старта &
	Возникает &
	Возникает &
	Не возникает при наличии начального профиля \\ \hline
	
	Проблема белой вороны &
	Возникает &
	Возникает &
	Не возникает \\ \hline
	
	Проблема устаревания пользовательских предпочтений &
	Возникает, если не учитывать новые оценки &
	Возникает, если не обновлять историю взаимодействий &
	Не возникает при своевременном обновлении предпочтений \\ \hline
	
\end{longtable}